\subsection{Perturbation of injective or surjective linear maps}

THEOREM 1. --- Let E and F be separated locally convex spaces and let u be a strict morphism from E into F, whose kernel is finite-dimensional and whose image is closed. Let h be a compact linear map from E into F. The linear map \(u + h\) is a strict morphism, its kernel is finite-dimensional, and its image is closed.

By prop. 1 of III, p. 71, there exists a closed neighborhood V of 0 in E such that the restriction of \(u\) to V is proper. Since \(h\) is compact, there exists a closed neighborhood W of 0 contained in V such that the set \(h(\mathrm{W})\) is relatively compact. Set \(\mathrm{C} = \overline{h(\mathrm{W})}\) . The restriction of \(u\) to W is proper (TG, I, p. 74, cor. 1). The map \(\alpha: x \mapsto (u(x), h(x))\) from W into \(\mathrm{F} \times \mathrm{C}\) is proper because the map \(\mathrm{pr}_1 \circ \alpha = u|\mathrm{W}\) from W into F is proper (TG, I, p. 73, prop. 5, \(d\) ). The map \(\beta: (y, z) \mapsto (y + z, z)\) from \(\mathrm{F} \times \mathrm{C}\) into \(\mathrm{F} \times \mathrm{C}\) is a homeomorphism, and the map \(\mathrm{pr}_1\) from \(\mathrm{F} \times \mathrm{C}\) into F is proper (TG, I, p. 77, cor. 5). The composite map \(\mathrm{pr}_1 \circ \beta \circ \alpha\) from W into F is therefore proper (TG, I, p. 73, prop. 5, \(a\) ). But this map is none other than the restriction of \(u + h\) to W. By prop. 1 of III, p. 71, \(u + h\) is a strict morphism, its kernel is finite-dimensional, and its image is closed.

Lemma 2. --- Let E and F be Fréchet spaces and \(u \in \mathcal{L}(\mathrm{E};\mathrm{F})\) . The following conditions are equivalent:

\begin{enumerate}
\def\labelenumi{(\roman{enumi})}
\item
  The cokernel of u is finite-dimensional;
\item
  \(^{t}u \colon \mathrm{F}_c' \to \mathrm{E}_c'\) is a strict morphism with closed image whose kernel is finite-dimensional.
\item
  \(\Longrightarrow\) (ii): Suppose that the cokernel of u is finite-dimensional. The map u is then a strict morphism (Lemma 6 of III, p. 52). By Prop. 2 of EVT, IV, p. 27, the image of \(^{t}u\) is closed in E' equipped with the weak topology, and a fortiori in \(E_{c}'\) . The kernel of \(^{t}u\) is the orthogonal complement of the image of u (loc. cit.); it is therefore finite-dimensional. Finally, \(^{t}u\) is a strict morphism from \(F_{c}'\) into \(E_{c}'\) by Theorem 1 of EVT, IV, p. 28.
\item
  \(\Longrightarrow\) (i): Suppose that \(^t u\colon F_c' \to E_c'\) is a strict morphism. By EVT, IV, p. 28, th. 1, the image of \(u\) is closed. The kernel of \(^t u\) is the orthogonal complement of \(\operatorname{Im}(u)\) (by EVT, IV, p. 27, prop. 2); if this kernel is finite-dimensional, the image of \(u\) is of finite codimension in F.
\end{enumerate}

THEOREM 2. --- Let E and F be Fréchet spaces, u: E → F a continuous linear map whose cokernel is finite-dimensional, and h: E → F a compact linear map. The linear map u + h is a strict morphism, its cokernel is finite-dimensional, and its image is closed.

By lemma 6 of III, p. 52, it suffices to show that the cokernel of \(u + h\) is of finite dimension. Now, by prop. 9 of III, p. 9, the map \(^{t}h\) from \(F_{c}^{\prime}\) into \(E_{c}^{\prime}\) is compact. Theorem 2 follows from Theorem 1 and Lemma 2.

